\chapter{Introduction}
\label{ch:intro}

A supernumerary robotic limb is an extra arm worn on the body, working
alongside the wearer's own limbs rather than replacing or rehabilitating
them~\cite{yang2021srlreview,prattichizzo2021augmentation}. The wearer keeps
full natural function, which separates these devices from prostheses. What
separates them from fixed manipulators is more consequential for control: the
base of the robot is a person. It accelerates, sways and steps. It is soft and
irreplaceable. It cannot be told to hold still.

This thesis reports work on the MUlti-limb Virtual Environment, known within
Imperial College London as the Dr.\ Octopus platform~\cite{muve}, which mounts
two Kinova Gen3 seven-degree-of-freedom arms~\cite{kinova2022gen3} on a
backpack frame. The problem addressed is the command interface. Two redundant
seven-degree-of-freedom arms present fourteen commandable joints, and a person
has two hands and a limited attention budget. Something has to close that gap.

\begin{figure}[htbp]
  \centering
  % Who is where, and what passes between them.
\begin{tikzpicture}[x=1mm, y=1mm, font=\scriptsize]

  \node[person, minimum height=22mm, text width=30mm, fill=violet!10] (drv)
        at (0,0) {\textbf{The driver}\\[2pt] stands across the room, moves a
        sensed mannequin arm, or holds two hand controllers, or types a
        sentence};

  \node[stage, minimum height=22mm, text width=32mm] (sw) at (48,0)
        {\textbf{The software}\\[2pt] turns the input into a wanted hand
        position, checks it against the wearer's body, and sends joint
        commands};

  \node[person, minimum height=22mm, text width=30mm, fill=violet!10] (wr)
        at (96,0) {\textbf{The wearer}\\[2pt] carries \SI{17}{\kilo\gram} on a
        backpack frame and does not control the arms};

  \node[motor, minimum height=10mm, text width=32mm] (arms) at (48,-26)
        {\textbf{Two seven-jointed arms}\\ mounted behind the wearer's
        shoulders};

  \draw[go] (drv) -- node[tag, above]{intent} (sw);
  \draw[go] (sw) -- (arms);
  \draw[go] (arms.east) -- ++(20,0) |- node[tag, right, pos=0.25]
        {motion beside the head} (wr.south);
  \draw[go, violet!60!black] (wr.north) -- ++(0,10) -| node[tag, above, pos=0.25]
        {where the body is} (sw.north);
  \draw[go, dashed] (arms.west) -- ++(-22,0) |- node[tag, left, pos=0.25]
        {what the cameras see} (drv.south);

  \node[guard, minimum height=8mm, text width=44mm] (stop) at (48,-42)
        {\textbf{A third person} holds the emergency stop and watches the
        wearer, not the screen};
  \draw[halt] (stop) -- (arms);

\end{tikzpicture}

  \caption[The two-person arrangement]{The arrangement this thesis addresses.
  The operator commands the arms and cannot feel them; the wearer carries them
  and does not command them. Neither person is in a position to guarantee the
  wearer's safety, so the software between them is, and that allocation is
  what the safety architecture of \cref{ch:method} exists to discharge.}
  \label{fig:overview}
\end{figure}

\section{Two design premises}

The work rests on two premises that were adopted early and shaped everything
downstream. Both are stated here because the rest of the thesis only makes
sense if they are visible from the start.

\subsection{Commanding fourteen joints is a usability problem, not a
plumbing problem}

The naive interface exposes every joint. It is also unusable. A body-worn
master that reproduces the slave's full kinematic chain requires the operator
to coordinate seven degrees of freedom per arm through finger motion, with no
force feedback and no view of the joint the arm is actually in. Earlier work
on miniature physical masters makes the case for morphological similarity
between master and slave~\cite{guggenheim2021haptic}, and that case is
accepted here. What does not follow is that every degree of freedom of the
slave must be commanded independently by the operator.

The architecture adopted instead splits the command by what each source
measures well. Position comes from three or four well-conditioned
potentiometer channels resolved into a spherical form about the master's
shoulder. Orientation comes from the inertial measurement unit at the wrist,
which observes gravity directly and therefore needs no integration and no
calibration against an absolute reference. What is left over, principally the
wrist orientation required to grasp a particular object, is supplied by
autonomy or is not commanded at all.

This is a reduced-degree-of-freedom design by intent. It is the reason the
system uses a spherical decomposition rather than forward kinematics from
seven joint angles, and the reason its default orientation mode pins the wrist
rather than tracking it. Both choices cost capability, and both are quantified
in \cref{ch:results}.

\subsection{A body-worn master will lose channels, so losing them must be
survivable}

The master carries fourteen analog channels routed through seven rotating
joints per arm. Connectors fatigue. Wires flex through joints that turn. Over
the life of any such device some fraction of its channels will stop reporting
usable values, and the fraction will grow. A design whose behaviour is
undefined when one channel fails is not a design that has been finished.

The system therefore treats channel loss as an operating condition rather than
a fault. Each channel carries a health verdict computed from its own recent
history. Validation is scoped to the channels the active position mode
consumes, so a dead channel that the mode does not read cannot reject a frame.
Individual channels can be disabled at run time, and the resulting loss of
capability is published rather than inferred. When enough channels are lost
that position can no longer be estimated, the system reports position
unavailable instead of extrapolating one.

That last behaviour is the design's centre of gravity. On a manipulator bolted
beside a person's head, a fabricated position is worse than no position, and
the difference between them has to be visible from outside the process.

\section{Objectives}

The planning report framed this project around a haptic master and slave pair
with physical springs at the master and virtual impedance at the slave. The
delivered system diverged from that plan, for reasons set out in
\cref{ch:discussion} and recorded in full in the accompanying divergence
analysis, which is retained as a supporting file rather than as an appendix
because it documents the route rather than the result. The objectives below
are the revised
set, aligned to what was built and measured.

\begin{enumerate}
  \item To design and build an instrumented mannequin master reproducing the
    Gen3 morphology at reduced scale, and to characterise its sensing against
    recorded data rather than against specification.
  \item To derive a mapping from master pose to slave end-effector command
    that degrades predictably as sensing channels are lost, and to quantify
    what each loss costs.
  \item To implement a safety architecture appropriate to a manipulator
    mounted beside a person, and to validate it by injecting faults rather
    than by inspecting code.
  \item To characterise the reachable workspace of the mounted configuration
    and determine which manipulation tasks it can support.
  \item To establish a verification methodology producing traceable evidence
    for every claim, and to apply it across the task set.
  \item To identify, by measurement, what must be resolved before hardware
    validation and a user study become possible.
\end{enumerate}

\section{Challenges}

Four difficulties consumed most of the effort.

\emph{The master and the slave are not the same shape.} The master reproduces
the Gen3's joint sequence but not its proportions. Its links sum to
\SI{0.272}{\metre} against the arm's \SI{0.902}{\metre} reach, a ratio of
roughly three to one, so joint-to-joint correspondence is unavailable. Working
in task space instead raises a second problem, which is that the two chains
reach different sets of poses. The resolution was to map displacement of the
master's tip onto displacement of the slave's end effector about an anchor
that can be re-latched, described in \cref{ch:method}.

\emph{Two coordinate conventions meet at the driver.} The manufacturer's
controller reports joint positions in degrees over $[0, 360)$; the middleware
uses radians over $(-\pi, \pi]$. The failure this invites is not a small
error. An angle of \SI{350}{\degree} read as $+350$ rather than $-10$ commands
very nearly a full revolution of a joint standing next to a person. The
conversion was isolated in one module and tested at the seams.

\emph{Calibration references must survive a power cycle without producing a
step.} The potentiometers report absolute values whose meaning depends on how
the master was assembled, so a zero reference has to be captured and stored.
If that reference is taken from a frame that failed validation, or if the
commanded pose is referenced to the origin of the master's own frame rather
than to where the operator's hand actually is, then the first command issued
after start-up is a step of the master's entire tip vector.

\emph{Almost every failure in this system is silent.} An arm that has stopped
moving looks the same whether it is blocked for a good reason, blocked for a
bad one, or receiving no data at all. A large part of the engineering effort
went into making internal state legible from outside. One of the more
uncomfortable findings of this work is that the measuring instruments failed
more often, and more quietly, than the system being measured.

\section{Contributions}

Each contribution is stated with its maturity, because the distinction between
implemented, verified in simulation, and validated on hardware is load-bearing
throughout. \textbf{No contribution in this thesis was validated on physical
hardware.} \Cref{ch:discussion} explains why.

\begin{enumerate}
  \item \textbf{A graceful-degradation architecture for a body-worn master.}
    Per-channel health classification computed over distinct sensor updates,
    validation scoped to the channels the active mode consumes, run-time
    channel disabling with published capability loss, and a reduced-degree-of-freedom
    fallback that keeps position available on a single potentiometer.
    The cost of each rung is measured on recorded hardware data.
    \emph{Implemented; measured against \num{20430} recorded frames.}

  \item \textbf{A reduced-degree-of-freedom command mapping} that takes
    direction from gravity and magnitude from the best-conditioned joints,
    with an anchor that re-latches on clutch engagement so indexing is
    unbounded and jump-free. \emph{Implemented; verified in simulation.}

  \item \textbf{A safety architecture whose blocking is observable.} Every
    mechanism able to stop the arm registers by name, must state how it
    clears, and expires into an explicit unknown state if the loop asserting
    it stops running. \emph{Implemented; fourteen of fourteen injected faults
    handled.}

  \item \textbf{A workspace characterisation of the mounted configuration},
    including the result that the two arms share no reachable workspace under
    the as-built mount. This blocks inter-arm handover for geometric reasons
    that no control or planning change can address. A buildable mount
    modification that would resolve it is quantified.
    \emph{Verified in simulation;} \cref{ch:workspace}.

  \item \textbf{A jerk-limited, phase-synchronised motion generator} replacing
    the per-joint step limiter that constituted the entire motion generation
    for teleoperation. It removes \SI{51.3}{\milli\metre} of unrequested
    Cartesian excursion on a platform whose grasp capture gate is
    \SI{30}{\milli\metre}, and brings peak commanded joint speed from
    \num{12.53} times the manufacturer's declared limit to \num{1.00}.
    \emph{Implemented; measured offline against the platform's own forward
    kinematics.}

  \item \textbf{A calibrate-then-plan stage} that measures the working surface
    and the objects on it instead of being given their coordinates, and hands
    the result to a planner. Six objects are recovered with every centre
    within \SI{1}{\milli\metre} and every extent within \SI{3}{\milli\metre},
    and the resulting map is the first thing in this project that a motion
    planner has ever consulted. \emph{Implemented; verified against the
    simulated scene and against mock depth.}

  \item \textbf{A positioning error budget} assembled from separately measured
    contributions, giving \SI{18.6}{\milli\metre} combined in quadrature and
    \SI{33.4}{\milli\metre} added, against a \SI{30}{\milli\metre}
    requirement, and identifying the camera mounting transform as the largest
    unmeasured term. \emph{Measured; one term explicitly unmeasured and
    labelled.}

  \item \textbf{A stage-by-stage instrumentation of the perception pipeline
    on four physical cameras.} Thirty-three stages per camera, each drawing
    itself on the frame it was given and reporting its quantities with units;
    \num{96} of \num{132} stages produced a result and none crashed. It
    establishes the surface fit on a real table at \SI{1.43}{\milli\metre}
    root mean square with \SI{74}{\percent} inlier support at working
    distance, separates a refusal from an absence by design, and found within
    the first hour a scene-camera orientation that three files had asserted
    and none had measured. \emph{Measured on physical cameras;}
    \cref{ch:vision}.

  \item \textbf{A validation of the study's own analysis against constructed
    ground truth.} The pre-registered procedure recovers the effect the design
    was powered for on \SI{97.3}{\percent} of \num{20000} simulated studies
    family-wise but on only \SI{63.8}{\percent} for any named task once the
    correction is applied, while rejecting a true null at exactly the nominal
    rate. The design's own power calculation reports neither number.
    \emph{Measured on constructed data;} \cref{sec:not-run},
    \cref{app:analysis}.

  \item \textbf{A reproducible simulation campaign.} Every simulation this
    project can run, behind one command, in two declared tiers, with each
    stage's status and duration recorded and the thesis's tables generated
    from that record rather than transcribed from a terminal. The generator
    refuses to emit a table whose stage did not run. \num{16} stages, all of
    which ran; \num{13927} inverse-kinematics calls at zero failures; and one
    stage that took four passes, because its own monotonicity control refused
    a result three times --- twice for genuine under-sampling, and once
    because the control was comparing a lower bound against a boundary. The
    control was corrected rather than the number, and every reach it reports
    is unchanged by the correction. \emph{Measured;} \cref{sec:campaign}.

  \item \textbf{A verification methodology} requiring an instrument to recover
    a known answer before a surprising measurement taken with it is reported.
    \emph{Applied throughout.}
\end{enumerate}

The last arose from necessity. Across this project, results that looked like
system failures turned out to be defects in the measuring tool on
twenty-six separate occasions, and validating the instrument first caught
errors that would otherwise have entered the record as findings. Two of them
would have removed results that are now central to this thesis. It is the most
transferable outcome of the work. \Cref{sec:validity} states what it means for
the measurements reported here, and \cref{app:instrument} is the whole body of
evidence rather than a selection.

\section{Report structure}

\Cref{ch:background} reviews the field by theme and ends with the gap this
work addresses.

\textbf{\Cref{ch:master,ch:method,ch:perception,ch:motion} are the methods:
what was designed and built.} \Cref{ch:master} documents the master arm --- its
mechanical chain, its electronics and its firmware. \Cref{ch:method}
describes the delivered system: the command path, the degradation
architecture, the six operating modes, the safety layer and the operating
window. \Cref{ch:perception} covers the perception stage that replaced
declared coordinates with measured ones, the planner that consults it, and
the natural-language route into the task layer. \Cref{ch:motion} covers
motion generation and what is known about transferring a joint command to a
physical arm, and assembles the positioning and reaction budgets.

\textbf{\Cref{ch:results,ch:vision,ch:workspace} are the results.}
\Cref{ch:results} reports the simulation campaign and every measurement made
of the system, each labelled with the strength of evidence behind it.
\Cref{ch:vision} reports the perception pipeline exercised on four physical
cameras, and is together with \cref{sec:hardware-session} the only part of
this report whose numbers come from hardware in the laboratory.
\Cref{ch:workspace} presents the workspace characterisation separately,
because it is a chain of measurements whose consequences determine the shape
of the task set.

\Cref{ch:discussion} interprets those results, states what the measurements
are worth (\cref{sec:validity}) and where they stop being trustworthy, and
\cref{ch:conclusion} states what was achieved against the objectives above.

The appendices carry what a final report should record but not argue from.
\Cref{app:development} is the route the design took, including the approaches
that were abandoned. \Cref{app:instrument} is the case-by-case evidence behind
the validity argument. \Cref{app:study,app:analysis} are the protocol for the
two-person study and the analysis that would evaluate it: the study was
designed in full and could not be run, for the reasons in
\cref{sec:not-run}.